\documentclass[10pt,a4paper]{amsart}
\usepackage[margin=19mm]{geometry}
\usepackage{amsmath,amssymb,mathtools}
\usepackage[scr=rsfs,cal=euler]{mathalfa}
\usepackage{xfrac}
\usepackage[unicode,hidelinks,bookmarks=false]{hyperref}
\newcommand{\ie}{i.e.\ }
\newcommand{\cf}{cf.\ }
\newcommand{\resp}{respectively\ }
\newcommand{\Subsection}{\subsection}
\providecommand{\nobreakdash}{\nobreak\hbox{-}}
\setlength{\emergencystretch}{2em}
\multlinegap0pt
\usepackage{amssymb}
\usepackage{enumerate}
\usepackage{bm}
\newtheorem{theorem}{Theorem}
\newtheorem{proposition}{Proposition}
\def \b {\beta}
\title{Vafa--Witten Equations and Conformal Geometry}
\author{Teng Huang, Pan Zhang}
\date{Source: arXiv:2606.22100v1 (CC BY 4.0)}
\begin{document}
\maketitle

\noindent\emph{Source and licence.} Vafa--Witten Equations and Conformal Geometry. Version arXiv:2606.22100v1 (CC BY 4.0), distributed by arXiv under the Creative Commons Attribution 4.0 International licence, http://creativecommons.org/licenses/by/4.0/, as recorded in the arXiv licence metadata for this exact version and shown on the abstract page https://arxiv.org/abs/2606.22100v1. Full licence notice: \texttt{licenses/arxiv-2606.22100-CC-BY-4.0.txt}.\par\smallskip

\noindent\textit{Editorial scope.} Counted unit: the article's theorem T3 (source0.tex lines 243--246). The article's complete original proof of it is delivered with it (source0.tex lines 689--712, one \texttt{proof} environment of 24 lines, which the article places away from the statement). No mathematics is written, repaired or re-derived: the statement and the proof are the article's own lines, in the article's own notation, and no retained line is substituted or re-flowed.  Retained uncounted as the source-local context the proof needs: lines 189--208; lines 227--229; lines 667--672.  Nothing was omitted from any retained span, so the package carries no omission record.  No retained line contains a citation, so the wrapper carries no bibliography and no bibliography entry.  No retained line contains a figure environment, an image-inclusion command or drawing code, so no image is delivered and the package declares no asset dependency.  Editorial formatting only, in addition: an AMS \texttt{amsart} preamble that restates the article's own declarations and macros used by the retained lines, namely the environments \texttt{theorem}, \texttt{proposition} on separate counters, together with the standard packages those lines need.  The retained environments print in this document as Theorem 1, Proposition 1 in order of appearance, whereas the article prints the same items with the numbering of its own full document.  The complete native source of the article as distributed in the arXiv e-print for this exact version is bundled unchanged as \texttt{sources/203324/source0.tex}.
\begin{theorem}\label{T1}
Let $(M,g)$ be a closed, smooth, four-dimensional manifold. Suppose $(A,B)$ is a solution to the Vafa-Witten equations with $B\not\equiv0$. 
Then the following hold:

\noindent\textup{(1)} The Yamabe constant $Y(g)$ satisfies the inequality
\begin{equation} \label{E0}
Y(g)\leq 2\sqrt{6}\|W^{+}_{g}\|_{L^{2}},
\end{equation}
where $W^{+}_{g}$ denotes the self-dual Weyl tensor of $g$.

\noindent\textup{(2)} If the Yamabe constant $Y(g)>0$, then the following topological bound holds:
\begin{equation}\label{E10}
\int_{M}|W^{+}_{g}|^{2}\,d\mathrm{vol}_{g}\geq\frac{4}{3}\pi^{2}\bigg{(}2\chi(M)+3\sigma(M)\bigg{)}.
\end{equation}

\noindent\textup{(3)} If $Y(g)=2\sqrt{6}\|W^{+}_{g}\|_{L^{2}}$, then 
\[
F_{A}^{+}=0\quad\text{and}\quad [B_{\bullet}B]=\nabla_{A}B=0.
\]
\end{theorem}

\begin{theorem}\label{T2}
Let $(M,g)$ be a closed, smooth, four-manifold. Suppose that $(A,B)$ is a solution to the Vafa--Witten equations with $B\not\equiv0$ and structure group $SU(2)$ or $SO(3)$. Then the equality $Y(g)=2\sqrt{6}\|W^{+}_{g}\|_{L^{2}}$ holds if and only if $(M,g)$ a K\"{a}hler manifold of non-negative scalar curvature and the connection $A$ is reducible.
\end{theorem}

\begin{proposition}\label{P5}
	Let $(M,g)$ be a closed K\"{a}hler surface with K\"{a}hler metric $g$. Let $(A,B)$ be a solution of the Vafa--Witten equations on $(M,g)$.
	 Suppose the scalar curvature of $g$ is positive. Then $(A,B)$ satisfies 
	$$F_{A}^{+}=0,\quad \b=0,\quad d_{A}\gamma=0.$$
	Moreover, if $A$ is an irreducible connection, then $B=0$. 
\end{proposition}

\begin{theorem}\label{T3}
Let $(M,g)$ be a closed, four-dimensional Einstein manifold of positive scalar curvature, normalized so that $Ric(g)=3g$. Suppose $(A,B)$ is a solution of Vafa-Witten equations on a principal $SU(2)$ or $SO(3)$-bundle $P\rightarrow M$ with $A$ is irreducible and $B\not\equiv0$. 
Then $(M,g)$ is not K\"{a}hler and satisfies \[\frac{2}{9}\pi^{2}\bigg{(}2\chi(M)+3\sigma(M)\bigg{)}>vol(g). \]
\end{theorem}

\begin{proof}[\textbf{Proof of Theorem \ref{T3}}]
Let $(M,g)$ be a closed, oriented four‑dimensional Einstein manifold with positive scalar curvature, normalized so that $\mathrm{Ric}(g)=3g$. For any such metric, the Chern–Gauss–Bonnet and signature formulas give
\begin{equation}\label{E14}
2\pi^{2}(2\chi(M)+3\sigma(M))=\int_{M}|W^{+}_{g}|^{2}dvol_{g}+\frac{1}{48}\int_{M}|R_{g}|^{2}dvol_{g}.
\end{equation}
By Theorem \ref{T1}, any Vafa–Witten solution $(A,B)$ with $B\not\equiv0$ satisfies
\begin{equation}\label{E15}
\int_{M}|W^{+}|^{2}\geq\frac{4}{3}\pi^{2}(\chi(M)+3\sigma(M)).
\end{equation}
Inserting (\ref{E15}) into (\ref{E14}) and using the normalization $R_{g}=12$, we obtain
\begin{equation*}
\begin{split}
2\pi^{2}(2\chi(M)+3\sigma(M))&\geq\frac{4}{3}\pi^{2}(\chi(M)+3\sigma(M))\\
&+\frac{1}{48}\int_{M}144dvol_{g}.\\
\end{split}
\end{equation*}
Consequently,
\begin{equation}\label{E9}
\frac{2}{9}\pi^{2}(2\chi(M)+3\sigma(M))\geq vol(g).
\end{equation}
Equality in (\ref{E9}) holds if and only if $Y(g)=2\sqrt{6}|W^{+}_{g}|_{L^{2}}$. By Theorem \ref{T2}, this would imply that $(M,g)$ is Kähler and that $A$ is irreducible, contradicting the hypothesis. Hence the inequality is strict. 

Moreover, if $(M,g)$ is K\"{a}hler, then by Proposition \ref{P5}, we would obtain  $B=0$, which again contradicts the assumption $B\not\equiv0$. Therefore, $M$ cannot be K\"{a}hler. This completes the proof of Theorem \ref{T3}.
\end{proof}

\end{document}
